\section{基于运行生产资料的汽轮机系统故障知识图谱构建}
\label{sec:knowledge_graph}

\subsection{引言}

上一章从正常运行数据出发，建立了面向燃煤机组汽轮机系统的设备级状态监测方法，其本质是在给定运行边界和操作状态下学习设备应具有的正常状态响应。然而，状态监测模型主要回答“当前运行状态是否偏离正常参考”，并不能仅依靠数值偏差完整回答“异常发生在哪个设备、哪些参数之间存在物理联系、异常可能沿何种路径传播以及对应何种故障机理”等诊断问题。对于汽轮机这类由多级汽缸、再热、抽汽回热、给水及冷端系统共同组成的复杂热力系统，不同设备之间存在明确的蒸汽流动、能量传递和控制关系，故障影响也可能沿这些关系向相邻设备和参数传播。因此，若仅将全部监测参数作为无结构向量输入分类模型，模型虽能够学习统计相关特征，却难以显式利用汽轮机系统已有的设备拓扑、热力机理和运行经验。

知识图谱（Knowledge Graph, KG）以实体、关系和属性的形式组织异构知识，能够将分散于系统结构、运行规程、设备资料、故障案例和学术文献中的非结构化知识转化为可计算的图结构\cite{Hogan2021KnowledgeGraphs}。对于工业故障诊断，知识图谱的价值并不在于替代实时运行数据，而在于提供“哪些参数属于哪个设备”“不同设备如何连接”“某类故障通常会引起哪些异常现象”“异常又可由哪些监测参数反映”等结构先验。已有研究已将过程知识与图神经网络结合，用于约束复杂工业系统的故障传播和根因分析\cite{Han2022KnowledgeEnhancedGNN}。针对汽轮机系统，Zhang和Hu根据设备机理构建过程图和故障传播图，并利用图神经网络推断汽封缺陷、蒸汽泄漏及汽耗增加之间的传播路径，验证了图结构描述汽轮机故障传播关系的可行性\cite{ZhangHu2021FaultPropagationGNN}。

基于上述认识，本章在杨昊东针对CFB锅炉运行生产资料构建领域知识图谱的技术路线基础上\cite{YangHaodong2026CFB}，面向燃煤机组汽轮机系统重新定义知识范围、实体模式和关系模式。首先结合汽轮机主通流、再热、抽汽回热、给水和冷端系统分析典型故障的传播特性；随后从系统拓扑、DCS测点、汽轮机运行与检修资料以及公开同行评议文献中组织故障知识，利用大语言模型生成候选三元组，并通过领域规则、语义对齐和证据校验完成知识筛选与实体标准化；最后构建“设备--参数--异常--故障--处置”五类实体组成的汽轮机故障知识图谱，为下一章KG-GCN模型提供静态参数关联拓扑。

\subsection{燃煤机组汽轮机系统参数及故障传播特性}
\label{subsec:turbine_fault_characteristics}

\subsubsection{汽轮机系统拓扑与故障传播路径}

汽轮机系统具有典型的层次化和网络化结构。主通流沿“主蒸汽--VHP--一次再热RH1--HP--二次再热RH2--IP--LP--凝汽器”的方向传递，抽汽则从不同压力级汽缸分流至高、低压加热器和除氧器，凝结水经低压加热、除氧、给水泵升压和高压加热后重新进入锅炉。与此同时，凝汽器通过循环水系统形成冷端边界。由此形成的设备网络同时包含主蒸汽流向、抽汽能量传递、给水回路和冷端换热等多类关系。

从故障传播角度看，汽轮机局部异常通常并不会只表现为单一测点变化。一方面，故障会改变故障设备内部正常的热力状态，例如汽封泄漏会改变蒸汽有效流量和效率，通流部件侵蚀或沉积会改变有效流通面积和级组压力分布；另一方面，这些状态变化还可能通过设备接口继续作用于下游系统。例如VHP排汽状态会影响一次再热入口，抽汽参数变化会作用于相应加热器，凝汽器背压变化则会反向影响LP末级膨胀过程。Zhang和Hu指出，大型工业系统中的故障可沿物质流或信息流路径从局部设备向系统边界传播，并在汽轮机算例中构建了由汽封缺陷、蒸汽泄漏到汽耗增加的分层故障传播图\cite{ZhangHu2021FaultPropagationGNN}。因此，将实际热力连接关系写入故障知识图谱，可以为后续模型区分“统计相关”与“具有工程意义的结构关联”提供先验约束。

为便于构建知识图谱，本文将汽轮机系统故障知识划分为三类关系。第一类为设备和参数的\textbf{静态结构关系}，包括设备归属、蒸汽流向、抽汽连接和冷端连接等；第二类为故障引起状态变化的\textbf{因果及表现关系}，例如“凝汽器水侧污垢--导致--换热能力下降”“换热能力下降--影响--凝汽器背压”；第三类为面向运行维护的\textbf{诊断与处置关系}，包括“故障--关联监测--参数”和“故障--处置为--检修动作”。上述关系共同构成从系统结构到故障机理再到运行处置的知识链。

\subsubsection{典型热力故障及参数响应特征}

汽轮机故障类型较多，既包括轴系、振动等机械故障，也包括通流、汽封、阀门、回热和冷端异常。本文研究重点为能够通过DCS热力参数反映并与前述状态监测模型衔接的热力/通流类故障，因此不将转子不平衡、轴承振动等机械振动故障作为本章知识图谱的主要对象。

（1）\textbf{调节阀死区及配汽异常。} 高压调节阀用于调节进入汽轮机的蒸汽流量并参与机组功率控制。由于阀杆、阀体及执行机构长期频繁动作，可能产生磨损、卡涩和死区。Zhang和Hu对汽轮机控制阀图模型的研究指出，死区表现为控制指令发生变化而实际阀位在一定范围内不发生响应，并可能诱发高压侧及系统振荡\cite{ZhangHu2021ControlValveGraph}。因此，在知识图谱中可形成“高压调节阀死区--导致--阀位对指令不敏感--表现为--阀位反馈异常”的故障链，并将主蒸汽流量和VHP入口压力波动作为相关监测信息。

（2）\textbf{汽封磨损与蒸汽泄漏。} 级间汽封和轴端汽封用于限制不同压力区域之间的蒸汽泄漏。汽封齿损伤、间隙增大或部件磨损会增加泄漏量。针对汽轮机故障传播的研究表明，HP、IP和LP的级间或轴端汽封缺陷可进一步导致相应部位蒸汽泄漏，并最终表现为为维持相同输出功率所需汽耗增加；对于LP轴端汽封，密封性能下降还可能造成空气漏入低压系统，影响凝汽器真空\cite{ZhangHu2021FaultPropagationGNN}。此外，针对HP--IP合缸中间汽封泄漏的热力研究表明，中间汽封泄漏是影响IP效率的重要因素，热再温度和IP排汽温度是泄漏量估计中的关键测量参数\cite{Xu2011HPIPLeakage}。因此，汽封故障能够形成由“部件缺陷--泄漏异常--效率/汽耗变化--相关温度和压力响应”组成的多层知识关系。

（3）\textbf{通流部件侵蚀、沉积及汽封磨损。} 长期运行过程中，蒸汽通流部件可能出现表面粗糙度增加、固体颗粒侵蚀、汽封磨损以及叶片表面沉积等退化。Kubiak等基于多台汽轮机现场诊断和检修结果指出，这些现象属于较常见的汽轮机通流退化形式；其中固体颗粒侵蚀可能造成首级喷嘴面积增大并改变关键压力分布，叶片沉积则会减小有效通流面积并导致容量和效率下降\cite{Kubiak2002TurbineDegradation,Kubiak2002ScaleDeposit}。这类故障的一个重要特点是：故障本身通常无法由单一DCS测点直接测量，但其对通流面积、压力分布、效率和出力形成的连锁影响可以由多个热力参数共同反映，因此特别适合采用“故障--异常现象--参数”分层方式进行图谱表达。

（4）\textbf{凝汽器污垢及冷端性能劣化。} 凝汽器水侧污垢会增加传热热阻并降低换热能力，是燃煤机组冷端常见性能劣化形式。Li等针对600~MW超临界燃煤机组提出凝汽器在线污垢监测模型，指出水侧沉积会降低凝汽器热有效性，并可通过污垢热阻变化实现在线监测\cite{Li2016CondenserFouling}。Medica-Viola等的煤电凝汽器运行状态模型进一步表明，凝汽压力受到汽轮机负荷、循环水流量、循环水入口温度及换热面污垢等因素共同影响；凝汽压力升高会减小汽轮机可利用的膨胀焓降并影响机组效率和出力\cite{Medica2018CondenserCondition}。因此，冷端知识图谱不仅需要包含“凝汽器污垢--换热性能下降--背压升高”的故障链，还应保留循环水温度和流量等外部边界，以避免将环境和操作变化误解释为设备故障。

（5）\textbf{给水加热器内部泄漏及回热性能下降。} 给水加热器内部管束、管板或隔板泄漏会改变正常的抽汽--给水换热过程。Barszcz和Czop针对煤电机组给水加热器建立用于模型诊断的动态模型，提出利用换热功率等过程指标跟踪慢性性能变化，包括由管道裂纹造成的内部泄漏\cite{Barszcz2011FeedwaterHeater}。Heo和Lee进一步研究了闭式给水加热器内部泄漏的定位与泄漏量诊断，指出该类泄漏会影响回热循环效率，并与相邻汽轮机和凝汽器状态形成耦合\cite{Heo2012FWHLeakage}。据此，可在图谱中建立“给水加热器内部泄漏--回热性能下降--机组热效率下降”的故障链，并将给水出口温度、抽汽侧压力等作为关联监测参数。

上述典型故障说明，汽轮机热力故障往往具有“故障原因不可直接观测、异常现象可由多个参数共同反映、影响可沿设备接口传播”的特点。若将原始DCS变量直接作为彼此独立的特征，难以显式表达这些结构关系；而知识图谱能够将设备拓扑、故障机理和监测参数组织为统一的图结构，为后续图卷积模型的邻域信息聚合提供物理依据。

\subsection{基于运行生产资料的知识图谱构建方法}
\label{subsec:kg_construction_method}

\subsubsection{汽轮机知识来源及证据组织}

知识图谱的可靠性取决于知识来源。杨昊东在CFB锅炉故障知识图谱中主要利用设计说明、运行日志和缺陷记录等运行生产资料，并通过大语言模型、NLP规则和语义对齐将自然语言知识转化为结构化三元组\cite{YangHaodong2026CFB}。面向汽轮机系统，本文将知识来源扩展为以下四类：

（1）\textbf{系统结构与DCS测点资料。} 包括汽轮机热力系统图、DCS画面、设备输入输出表和总测点表，用于确定设备、参数实体及“归属”“介质流向”“设备连接”等高可信结构关系。

（2）\textbf{运行规程、设备说明和检修资料。} 用于补充设备功能、正常运行边界、异常现象和运行处置知识。该类资料与具体机组结构高度相关，后续应随着现场资料补充持续更新。

（3）\textbf{汽轮机故障与性能诊断同行评议文献。} 对于现场历史故障记录暂不完备的故障类型，本文优先利用公开原始论文提取具有明确工程依据的“故障--现象--参数”关系。例如汽封泄漏、通流侵蚀、叶片沉积、凝汽器污垢和给水加热器内漏等关系均由相应原始研究支撑\cite{Kubiak2002TurbineDegradation,Xu2011HPIPLeakage,Li2016CondenserFouling,Heo2012FWHLeakage}。

（4）\textbf{既有故障诊断知识库及专家经验。} 该类资料用于补充异常描述同义词和处置动作，但不在缺少交叉验证时直接作为高置信因果边。早期汽轮机智能诊断研究已表明，专家系统可通过技术资料和运行人员经验组织大量故障现象、原因及处理知识，但其质量依赖知识获取和维护过程，因此需要进一步标准化和证据校验。

为了区分不同证据强度，本文在知识工程阶段为每条候选关系保留来源标识和证据等级。来自实际系统拓扑、真实DCS点表和直接支持该故障机理的原始同行评议文献的关系优先进入核心图谱；仅由相邻研究、经验推断或间接证据支持的关系暂作为候选边，后续结合运行资料或专家确认再决定是否纳入。通过保留“知识--证据”映射，可以避免图谱构建过程中将模型生成内容与真实工程知识混为一体。

\subsubsection{实体类型与关系模式定义}

知识图谱可表示为
\begin{equation}
    \mathcal{G}_{KG}=(\mathcal{V},\mathcal{E},\mathcal{R}),
    \label{eq:kg_definition}
\end{equation}
其中，$\mathcal{V}$表示实体集合，$\mathcal{E}$表示实体之间的关系边，$\mathcal{R}$表示关系类型集合。借鉴杨昊东知识图谱的实体组织方式，并结合汽轮机故障诊断需求，本文将实体集合划分为设备实体、参数实体、异常现象实体、故障实体和处置动作实体：
\begin{equation}
\mathcal{V}=\mathcal{V}_{D}\cup\mathcal{V}_{P}\cup\mathcal{V}_{A}\cup\mathcal{V}_{F}\cup\mathcal{V}_{O}.
\end{equation}
其中，$\mathcal{V}_{D}$包括VHP、HP、IP、LP、再热系统、凝汽器、加热器和阀门等设备或功能单元；$\mathcal{V}_{P}$表示可测或可计算的压力、温度、流量、阀位、效率等参数；$\mathcal{V}_{A}$表示“蒸汽泄漏”“真空下降”“通流面积减小”等状态异常；$\mathcal{V}_{F}$表示具体故障或性能退化原因；$\mathcal{V}_{O}$表示运行或检修处置动作。

关系模式根据汽轮机知识的物理语义定义为：\textbf{归属}，用于建立参数与设备的隶属关系；\textbf{介质流向/介质能量连接}，用于描述设备之间的蒸汽、水或能量传递；\textbf{发生于}，表示故障所在设备；\textbf{导致}和\textbf{影响}，用于表达故障或异常之间的因果作用；\textbf{表现为}，用于连接抽象异常与可观测状态；\textbf{关联监测}，用于连接故障与具有诊断价值的参数；\textbf{处置为}，用于连接故障与运行检修动作。通过限定实体类型和关系的合法组合，可以减少自然语言抽取产生的语义混乱。例如“VHP排汽温度--归属--VHP”为合法结构关系，而“VHP排汽温度--处置为--检查汽封”则不符合模式约束，应在规则校验阶段被剔除。

\begin{table}[htbp]
\centering
\caption{汽轮机故障知识图谱实体与关系模式}
\label{tab:kg_schema}
\begin{tabular}{p{2.2cm}p{5.2cm}p{6.2cm}}
\toprule
类别 & 定义 & 典型示例 \\
\midrule
设备实体 & 汽轮机设备、部件或功能单元 & VHP、RH1、HP、凝汽器、高压调节阀、给水加热器 \\
参数实体 & 可测或可计算的运行参数 & 主蒸汽压力、排汽温度、背压、阀位、循环水流量 \\
异常实体 & 对偏离正常状态的描述 & 蒸汽泄漏、真空下降、通流面积减小、系统振荡 \\
故障实体 & 造成异常的故障或退化模式 & 汽封缺陷、固体颗粒侵蚀、调节阀死区、凝汽器污垢 \\
处置实体 & 运行调整或检修动作 & 汽封检查、通流部件清洗、凝汽器水侧清洗 \\
\midrule
主要关系 & 设备、参数、异常、故障和处置之间的语义边 & 归属、介质流向、发生于、导致、影响、表现为、关联监测、处置为 \\
\bottomrule
\end{tabular}
\end{table}

\subsubsection{候选三元组抽取与领域规则校验}

运行规程、故障案例和学术文献通常采用自然语言描述，同一故障可能跨越多个句子才能形成完整的因果链。例如“汽封磨损增大泄漏量，导致单位功率汽耗上升”同时包含故障实体、异常实体和因果关系。为降低纯人工整理成本，本文采用大语言模型对文本进行候选三元组抽取，将知识统一转换为
\begin{equation}
    \tau_k=(h_k,r_k,t_k,s_k),
\end{equation}
其中$h_k$和$t_k$分别为头实体和尾实体，$r_k$为关系类型，$s_k$记录原始证据来源。大语言模型在该阶段只承担候选知识生成作用，其输出不能直接视为最终工程知识。

候选三元组随后经过领域规则校验。规则主要包括：（1）实体类型约束，例如“发生于”的尾实体必须为设备或部件；（2）关系方向约束，例如“参数--归属--设备”与“设备--介质流向--设备”的方向必须统一；（3）测点存在性约束，涉及本机组DCS参数时需能够在标准点表中找到对应实体；（4）物理路径约束，对设备传播边优先检查是否与实际主通流、抽汽、给水和冷端拓扑一致；（5）证据约束，故障因果关系需能够回溯至原始文本或工程资料。通过上述规则，可以将语言模型的开放式抽取结果压缩到受汽轮机领域模式约束的候选知识空间。

\subsubsection{实体标准化与关系置信度}

运行生产资料中常存在同义词、缩写和不同命名方式，例如“凝汽器压力”“凝汽器背压”和“排汽背压”在特定语境下可能指向相近概念，“调节阀反馈”“阀位反馈”和具体DCS点名也可能描述同一类参数。如果直接保留所有文本表述，会造成图谱节点重复并降低后续参数映射的一致性。因此，本文在规则校验后进一步进行实体标准化。

对于候选实体$e_i$和标准实体$e_j$，分别计算字符串相似度$S_{str}$和语义向量相似度$S_{sem}$。语义表示可利用Sentence-BERT等句向量模型获得\cite{Reimers2019SentenceBERT}，综合相似度表示为
\begin{equation}
S(e_i,e_j)=\lambda S_{sem}(e_i,e_j)+(1-\lambda)S_{str}(e_i,e_j),
\label{eq:entity_similarity}
\end{equation}
其中$\lambda$为语义相似度权重。当综合相似度超过阈值且实体类型一致时，将候选实体归并到标准词表；对于无法自动判断的设备专有名词和DCS测点，保留人工核验环节。本文不预先固定$\lambda$和阈值，而是在完成汽轮机标准词表和现场文本样本后再进行标定，以避免给出缺少实验依据的经验参数。

关系可靠性进一步由多源证据共同评价。借鉴课题组已有知识图谱构建思路，可将第$k$条关系的置信度表示为
\begin{equation}
C_k=w_1C_{LLM,k}+w_2C_{rule,k}+w_3C_{sem,k}+w_4C_{freq,k},
\quad \sum_{i=1}^{4}w_i=1,
\label{eq:relation_confidence}
\end{equation}
其中$C_{LLM,k}$表示语言模型对关系的抽取置信信息，$C_{rule,k}$表示领域规则校验结果，$C_{sem,k}$表示实体标准化与上下文语义一致性，$C_{freq,k}$表示该关系在独立资料或文本证据中的重复支持程度。需要指出的是，式~\eqref{eq:relation_confidence}用于定义后续知识筛选框架，当前核心图谱优先采用可由系统拓扑或原始文献直接支持的关系，尚未将未经标定的权重组合用于生成最终定量置信度。

\subsection{汽轮机系统故障知识图谱构建结果}
\label{subsec:kg_result}

\subsubsection{核心语义知识图谱}

依据上述模式，本文首先构建汽轮机故障知识图谱的核心语义层。当前初始版本以已确认的汽轮机系统拓扑、DCS参数类别和公开原始故障研究为主要证据，覆盖主通流、再热、抽汽回热、给水和冷端系统，并聚焦可由热力运行参数反映的通流/性能故障。经实体去重与关系整理后，初始图谱包含107个语义节点和114条关系三元组，其中设备实体23个、参数实体39个、异常实体24个、故障实体15个、处置实体6个。关系中包含36条参数归属关系、22条以上故障/异常因果关系，以及设备介质连接、故障发生位置、关联监测和处置等关系。上述规模仅表示当前用于论文方法开发的\textbf{核心语义图谱}，后续随着运行规程、缺陷记录及全系统DCS测点映射补充，图谱节点和边仍将继续扩展，因此不将该规模作为最终工程知识库规模。

\begin{table}[htbp]
\centering
\caption{汽轮机故障核心语义知识图谱初始规模}
\label{tab:kg_stats}
\begin{tabular}{lccccc}
\toprule
实体类型 & 设备 & 参数 & 异常 & 故障 & 处置 \\
\midrule
节点数 & 23 & 39 & 24 & 15 & 6 \\
\bottomrule
\end{tabular}
\end{table}

代表性知识链包括：
\begin{equation}
\text{HP级间汽封缺陷}\rightarrow\text{蒸汽泄漏}\rightarrow\text{汽耗率升高},
\end{equation}
\begin{equation}
\text{凝汽器水侧污垢}\rightarrow\text{换热能力下降}\rightarrow\text{凝汽器背压变化}\rightarrow\text{LP有效焓降下降},
\end{equation}
以及
\begin{equation}
\text{高压调节阀死区}\rightarrow\text{阀位对指令不敏感}\rightarrow\text{蒸汽流量/压力波动}.
\end{equation}
这些知识链并非单纯由文本相似性生成，而分别能够回溯到汽封故障传播、凝汽器状态监测和控制阀故障的原始研究\cite{ZhangHu2021FaultPropagationGNN,Medica2018CondenserCondition,ZhangHu2021ControlValveGraph}。通过将多类故障链与实际设备拓扑叠加，知识图谱能够同时表达“故障发生在哪里”“会产生什么异常”“哪些参数能够反映这种异常”“异常可能向哪些设备接口传播”。

\subsubsection{领域参数节点与DCS测点映射}

核心语义图谱中的“VHP排汽温度”“一级抽汽压力”等参数实体属于领域概念，而实际机组DCS中可能存在A/B侧、多冗余传感器以及具体KKS点号。为了将知识图谱用于后续KG-GCN，需要进一步建立\textbf{领域参数节点--实际测点节点}之间的映射层。设领域参数实体为$p_i$，对应实际DCS测点集合为
\begin{equation}
\mathcal{M}(p_i)=\{m_{i1},m_{i2},\ldots,m_{in_i}\},
\end{equation}
则模型可以将不同冗余测点的动态状态特征映射至同一领域语义位置，同时保留具体传感器节点用于后续图计算。

以第二章VHP为例，模型真实预测27个候选输出，最终保留24个状态参数作为后续诊断节点。第三章知识图谱并不要求将这24个测点重新定义为24种独立故障知识，而是先将其映射到VHP本体状态、一级抽汽接口、VHP--RH1冷再接口以及缸体/内部蒸汽等领域参数概念，再依据“参数--归属--设备”“设备--介质流向--设备”及“故障--关联监测--参数”等关系形成参数子图。对于后续HP、IP、LP、再热及凝汽器等模型，同样采用该方式将实际测点挂接到统一汽轮机语义图谱，从而实现不同设备模型与同一故障知识空间之间的连接。

这种“语义层+测量层”的两层设计可以避免直接以DCS点号作为知识图谱全部语义节点所带来的可读性问题：领域知识层保留稳定的汽轮机设备和故障语义，测量层负责适配具体机组的传感器配置。下一章将在此基础上进一步提取与状态监测参数对应的图邻接关系，并将状态监测模型产生的动态偏差信息映射至参数节点，构成KG-GCN的图结构输入。

\subsection{本章小结}

本章针对燃煤机组汽轮机系统故障知识分散于系统拓扑、DCS测点、运行资料和故障文献、难以直接用于数据驱动诊断的问题，构建了面向KG-GCN的汽轮机故障知识图谱方法。首先从主通流、再热、抽汽回热、给水和冷端系统出发，分析了调节阀死区、汽封泄漏、通流侵蚀/沉积、凝汽器污垢以及给水加热器内漏等典型故障的传播关系和参数响应；随后定义设备、参数、异常、故障和处置五类实体，以及归属、介质连接、导致、影响、表现、关联监测和处置等关系类型；在此基础上，采用“LLM候选抽取--领域规则校验--语义对齐--证据评价”的流程将非结构化知识转换为结构化三元组。

当前已形成由107个核心语义节点和114条证据三元组构成的汽轮机故障知识图谱初始版本。该图谱规模并非最终工程知识库规模，而是后续算法开发所需的核心语义层；随着全系统DCS测点、运行规程和现场故障资料继续补充，可进一步扩展设备和测量节点。下一章将在该知识图谱基础上构造与状态监测参数相对应的图拓扑，并将设备健康模型输出与实际运行状态之间的动态偏差信息映射至图节点，形成基于KG-GCN的汽轮机系统故障诊断模型。
